\documentclass[10pt,letterpaper]{article}

\usepackage[margin=0.52in]{geometry}
\usepackage[T1]{fontenc}
\usepackage[utf8]{inputenc}
\usepackage[default,type1]{sourcesanspro}
\usepackage[final]{microtype}
\usepackage[explicit]{titlesec}
\usepackage{enumitem}
\usepackage{tabularx}
\usepackage{array}
\usepackage{ragged2e}
\usepackage{iftex}
\usepackage[hidelinks]{hyperref}

\ifPDFTeX
  \input{glyphtounicode}
  \pdfgentounicode=1
\fi

\hypersetup{
  pdftitle={Connor Savugot Resume},
  pdfauthor={Connor Savugot},
  pdfsubject={Embedded Firmware and Embedded Linux Engineering},
  pdfkeywords={embedded firmware, C, C++, Embedded Linux, Yocto, Arago, Linux kernel, device tree, systemd, Docker, DSP, PIC32, ST microcontroller, MSP430FR2355, CAN, firmware update, power conversion, BMS}
}

\pagestyle{empty}
\setcounter{secnumdepth}{0}
\setlength{\parindent}{0pt}
\setlength{\parskip}{0pt}
\hfuzz=0.4pt

\titleformat{\section}
  {\fontsize{11.3}{11.9}\selectfont\bfseries\uppercase}
  {}{0pt}{#1\hspace{0.12cm}\titlerule[0.65pt]}
\titlespacing{\section}{0pt}{0.15cm}{0.055cm}

\newlist{resumebullets}{itemize}{1}
\setlist[resumebullets]{
  label=\textbullet,
  leftmargin=0.42cm,
  topsep=0.005cm,
  itemsep=0.005cm,
  parsep=0pt,
  partopsep=0pt
}

\newcommand{\companyheader}[1]{%
  {\fontsize{10.4}{10.8}\selectfont\bfseries #1}\par\vspace{-0.03cm}
}

\newcommand{\resumeentry}[2]{%
  \begin{tabularx}{\textwidth}{@{}X>{\raggedleft\arraybackslash}p{2.40in}@{}}
    \textbf{#1} & #2
  \end{tabularx}\vspace{-0.075cm}
}

\newcommand{\roleentry}[2]{%
  \begin{tabularx}{\textwidth}{@{}X>{\raggedleft\arraybackslash}p{2.40in}@{}}
    \textit{#1} & #2
  \end{tabularx}\vspace{-0.075cm}
}

\begin{document}
\fontsize{9.8}{10.72}\selectfont
\setlength{\RaggedRightRightskip}{0pt plus 5em}
\RaggedRight
\emergencystretch=1em
\hyphenpenalty=10000
\exhyphenpenalty=10000

\begin{center}
  {\fontsize{21.5}{22.2}\selectfont\bfseries Connor Savugot}\\[1.5pt]
  {\bfseries B.S. in Computer Engineering | North Carolina State University | Graduated May 2026}\\[1.5pt]
  {\fontsize{9.6}{10.2}\selectfont
    \href{mailto:consav04@gmail.com}{consav04@gmail.com}
    \enspace|\enspace \href{tel:+18285419487}{+1-828-541-9487}
    \enspace|\enspace Murphy, NC
    \enspace|\enspace \href{https://linkedin.com/in/connorsavugot}{linkedin.com/in/connorsavugot}
    \enspace|\enspace \href{https://github.com/Clem085}{github.com/Clem085}}
\end{center}
\vspace{-0.18cm}

\section{Technical Profile}
Embedded firmware engineer with experience in multi-MCU firmware boot processes, TI Embedded Linux platforms, programmable power-conversion control, and CAN-connected battery systems. Experienced across C/C++, multi-chip communication solutions, Linux services, board bring-up, bench debugging, and cross-platform Windows/Linux tooling.

\section{Technical Skills}
\textbf{Firmware \& OS:} Embedded C/C++, Python, Bash; TI Embedded Linux, Yocto/Arago, Linux kernel/device trees, systemd, Docker, FreeRTOS\\
\textbf{Processors \& Interfaces:} TI AM62, TI TMS320 DSP, PIC32, PIC33, STM32G474RE, MSP430FR2355; CAN J1939, I2C, SPI, UART, RS-232, ADC, PWM, GPIO, IPMI/IPMB\\
\textbf{Development \& Test:} Git/GitLab, Code Composer Studio, JTAG, KiCad, SocketCAN/PCAN; oscilloscopes, logic/protocol analyzers, through-hole/SMT soldering

\section{Engineering Experience}
\companyheader{AEGIS Power Systems}
\roleentry{Embedded Firmware Engineer}{May 2026--Present}
\begin{resumebullets}
  \item Build and debug TI AM62-based Embedded Linux systems using Yocto/Arago images and toolchains; compile and integrate kernels and device trees, diagnose boot and service failures, and develop systemd services for external-watchdog monitoring and recovery.
  \item Integrate and debug APU/EPU embedded power subsystems spanning battery hardware, CAN-connected devices, and Linux services; trace faults across power and wiring, bus traffic, Linux configuration, service logs, and application code.
  \item Diagnose CAN/CAN FD, IPMI/IPMB, I2C, SPI, and UART behavior with logs, register reads, analyzers, scopes, schematics, datasheets, and controlled substitutions to isolate hardware, firmware, timing, and configuration faults.
  \item Develop C/C++, Python, and Bash software, use Docker in development workflows, and improve code and tooling portability across Windows and Linux.
\end{resumebullets}

\roleentry{Embedded Systems Intern}{May--Aug 2025}
\roleentry{Computer Science Intern}{May--Aug 2024}
\roleentry{Selected engineering across both internships}{}
\begin{resumebullets}
  \item Developed TMS320 DSP firmware controlling transformer-coupled bidirectional buck--boost power stages through 12 PWM A/B pairs (24 physical outputs): independently configured each A output's period/frequency, duty cycle, and relative phase; paired B matched or inverted A.
  \item Developed multiple connected stages of one PIC32/TMS320 product's firmware-update workflow: a user-facing C++ CLI accepted supplied HEX files, Python/TCP transport provided chunking and checksum validation, and the processors coordinated external SPI-flash/golden-image handling, UART DSP programming, verification, and boot control.
  \item In 2024, built a restricted SSH/chroot CLI with Python/Bash for TI Yocto/Arago embedded power hardware; reverse-engineered RS-232 transactions and checksums; supported CAN/FreeRTOS debugging, schematic review, soldering, and burn-in-test hardware.
\end{resumebullets}

\resumeentry{TEAM Industries | Engineering Intern}{Jun--Aug 2021 \& Jun--Aug 2022}
\begin{resumebullets}
  \item Developed a repeatable Excel/VBA dimensional and spline-tolerance validation test that automatically flagged out-of-spec parts; built a searchable 17,000+ record tooling/part/vendor catalog with operation, location, and revision traceability.
\end{resumebullets}

\section{Senior Design}
\resumeentry{EV Active Sensor Adapter | ST MCU, TI BQ ADC, CAN}{NC State University}
\begin{resumebullets}
  \item Defined architecture and interfaces linking an STMicroelectronics MCU control subsystem, TI BQ-family battery-monitor/ADC, external CAN transceivers, and custom-PCB power, sensing, and telemetry hardware.
  \item Evaluated component and electrical/protocol tradeoffs; coordinated requirements, data paths, risk tracking, subsystem bring-up and validation, design reviews, and team documentation.
\end{resumebullets}

\section{Teaching \& Technical Leadership}
\resumeentry{ECE 306 Lead TA | NC State University}{Spring 2025--Spring 2026}
\begin{resumebullets}
  \item Supported three semesters: two in an unofficial role at about 10 hr/week and one as the official TA at 20 hr/week; delivered test preparation, labs, code/PCB debugging, project guidance, and ECE-program support.
  \item Taught and debugged vehicle subsystems spanning MSP430FR2355 firmware, ADC line sensors, timer/PWM motor control, GPIO/interrupt inputs, Wi-Fi/serial links, LCDs, and custom PCBs; added supplemental lectures and question-driven support as the class average rose from its typical C+ to A-.
\end{resumebullets}

\resumeentry{IEEE Open Project Space 2 | Lead Instructor}{Fall 2025 \& Spring 2026}
\begin{resumebullets}
  \item Ran two semesters of a project-based electronics course from idea through demonstration; taught firmware, PCB/BOM workflows, and through-hole/SMT assembly while coaching students to isolate and solve their own design problems.
\end{resumebullets}

\resumeentry{IEEE Soldering Workshops | Coordinator, Lecturer \& Lab Instructor}{}
\begin{resumebullets}
  \item Coordinated, lectured, and supervised standalone workshops using custom through-hole and surface-mount flashing-light PCB keepsakes; taught professional technique, safety, inspection, debugging, and rework.
\end{resumebullets}

\end{document}
